\section{随机变量究竟以什么方式接近}

同一列随机变量可以在某种意义下接近极限，在另一种意义下却不能。陈述极限定理之前，必须先指定“接近”是逐条样本路径、偏差概率、平均误差，还是仅指分布形状。

\begin{definition}[四种收敛]
在共同概率空间上的 \(X_n,X\) 满足：
\(X_n\to X\) 几乎处处，指
\(P(\{\omega:X_n(\omega)\to X(\omega)\})=1\)；
\(X_n\to X\) 依概率，指对每个 \(\varepsilon>0\)，
\(P(|X_n-X|>\varepsilon)\to0\)；
\(X_n\to X\) 在 \(L^p\) 中，指 \(p\ge1\) 且
\(\mathbb E|X_n-X|^p\to0\)。
依分布收敛 \(X_n\Rightarrow X\) 指对极限分布函数 \(F_X\) 的每个连续点 \(t\)，有 \(F_{X_n}(t)\to F_X(t)\)。最后一种定义只比较分布，\(X_n,X\) 可以定义在不同概率空间。
\end{definition}

若 \(X_n=X+1/n\)，四种收敛都成立。若每个 \(X_n\) 与 \(X\) 是独立的 Bernoulli\((1/2)\) 变量，则 \(X_n\) 与 \(X\) 同分布，因而 \(X_n\Rightarrow X\)，但
\(P(|X_n-X|>1/2)=1/2\)，所以不依概率收敛。这说明相同的分布不等于同一次实验的结果逐渐靠近。

\begin{proposition}[收敛方式的蕴含]
\(L^p\) 收敛或几乎处处收敛均蕴含依概率收敛；依概率收敛蕴含依分布收敛。一般不存在反向蕴含，但若 \(X_n\Rightarrow c\) 且 \(c\) 为常数，则 \(X_n\to c\) 依概率。
\end{proposition}
\begin{proof}
Markov 不等式给
\(P(|X_n-X|>\varepsilon)\le\mathbb E|X_n-X|^p/\varepsilon^p\)，证明第一条。
几乎处处收敛时，固定 \(\varepsilon>0\)，令
\(A_n=\{|X_n-X|>\varepsilon\}\)。事件
\(B_N=\bigcup_{n\ge N}A_n\) 随 \(N\) 递减，而
\(\bigcap_NB_N\) 是偏差超过 \(\varepsilon\) 无穷多次的事件，概率为零。概率测度的从上连续性给
\(P(A_n)\le P(B_n)\to0\)。若依概率收敛，则对每个有界连续 \(f\)，先把 \(|f(X_n)-f(X)|\) 分成 \(|X_n-X|\le\delta\) 与补集，在有界区间用一致连续性控制第一部分，再用尾概率控制区间外，得到 \(\mathbb E f(X_n)\to\mathbb E f(X)\)，故依分布收敛。最后，对 \(\varepsilon>0\)，分布函数在 \(c\pm\varepsilon\) 连续，且
\(P(|X_n-c|>\varepsilon)\le F_{X_n}(c-\varepsilon)+1-F_{X_n}(c+\varepsilon)\to0\)。
\end{proof}

反例还应区分几乎处处与依概率。在 \([0,1)\) 上按二进制区间逐层排列指标函数：第 \(m\) 层依次取长度 \(2^{-m}\) 的 \(2^m\) 个半开区间的指标。整列变量的非零概率趋零，故依概率收敛到零；但每条样本路径每层恰有一次为 \(1\)，不能逐点收敛。依概率收敛也未必给 \(L^1\) 收敛：令 \(X_n=n\) 的概率为 \(1/n\)，其余为零，则 \(X_n\to0\) 依概率，而 \(\mathbb E|X_n|=1\)。

\begin{theorem}[弱大数律，有限方差版]
若 \(X_1,X_2,\ldots\) 独立同分布，\(\mathbb E X_i=\mu\) 且
\(\operatorname{Var}(X_i)=\sigma^2<\infty\)，则
\(\overline X_n=n^{-1}\sum_iX_i\to\mu\) 依概率。
\end{theorem}
\begin{proof}
独立性消去协方差项，故
\(\operatorname{Var}(\overline X_n)=\sigma^2/n\)。
Chebyshev 不等式给
\(P(|\overline X_n-\mu|>\varepsilon)\le\sigma^2/(n\varepsilon^2)\to0\)。
有限方差只在这一简洁证明中使用；弱大数律另有更弱条件的版本。
\end{proof}

弱大数律只说每个固定 \(n\) 的偏差概率趋零，不保证在同一条无限
样本路径上从某处起始终接近均值。要升级为几乎处处收敛，
需要控制所有部分和中的最大偏差。

\begin{proposition}[Kolmogorov 最大不等式]
若 \(Y_1,\ldots,Y_n\) 独立、均值为零、方差有限，记
\(S_k=\sum_{i=1}^kY_i\)，则对任意 \(\lambda>0\)，
\[
P\!\left(\max_{1\le k\le n}|S_k|\ge\lambda\right)
\le\frac{\operatorname{Var}(S_n)}{\lambda^2}.
\]
\end{proposition}
\begin{proof}
令 \(A_k\) 是第 \(k\) 次首次达到 \(|S_k|\ge\lambda\) 的事件。
这些事件互不相交，且 \(A_k\) 只依赖前 \(k\) 个变量。
在 \(A_k\) 上写 \(S_n=S_k+(S_n-S_k)\) 并平方。
后续和与 \(A_k,S_k\) 独立且均值为零，所以交叉项的期望为零；
剩余的后续平方项非负。因此
\[
\mathbb E S_n^2
\ge\sum_{k=1}^n\mathbb E[S_k^2\mathbf1_{A_k}]
\ge\lambda^2\sum_{k=1}^nP(A_k).
\]
末项之和正是最大偏差事件的概率，且
\(\mathbb E S_n^2=\operatorname{Var}(S_n)\)，故得不等式。
\end{proof}

\begin{theorem}[强大数律，有限方差版]
若 \(X_i\) 独立同分布，\(\mathbb E X_i=\mu\) 且
\(\operatorname{Var}(X_i)=\sigma^2<\infty\)，则
\(\overline X_n\to\mu\) 几乎处处。
\end{theorem}
\begin{proof}
令 \(Y_i=X_i-\mu\)、\(S_k=\sum_{i=1}^kY_i\)。
对 \(n=2^m\) 及 \(\lambda=\varepsilon2^m\) 用最大不等式，
\[
P\!\left(\max_{k\le2^m}|S_k|\ge\varepsilon2^m\right)
\le\frac{\sigma^2}{\varepsilon^2\,2^m}.
\]
右侧对 \(m\) 求和有限。所用的 Borel--Cantelli 事实可在此直接证明：
若 \(\sum_mP(A_m)<\infty\)，则
\(P(\bigcup_{m\ge M}A_m)\le\sum_{m\ge M}P(A_m)\to0\)，
所以无限多个 \(A_m\) 同时发生的概率为零。
因此对每个正有理 \(\varepsilon\)，几乎每条路径从某个
\(m\) 起都满足
\(\max_{k\le2^m}|S_k|<\varepsilon2^m\)。
若 \(2^{m-1}<n\le2^m\)，便有
\(|S_n|/n\le2\max_{k\le2^m}|S_k|/2^m<2\varepsilon\)。
令有理 \(\varepsilon\downarrow0\)，得 \(S_n/n\to0\)
几乎处处，亦即所述结论。
\end{proof}

前两个大数律给平均值的极限，却没有给出有限样本误差的形状。若想把独立随机变量的和变成乘积，可以对分布作 Fourier 变换。实随机变量 \(Y\) 的\emph{特征函数}定义为
\(\varphi_Y(t)=\mathbb E\ee^{\ii tY}\)，\(t\in\mathbb R\)。它总存在，因为被平均的复数模长为一；独立性使
\(\varphi_{Y_1+\cdots+Y_n}(t)=\prod_j\varphi_{Y_j}(t)\)。这正是它比直接处理和的密度更方便的地方。
密度
\(f(z)=(2\pi)^{-1/2}\ee^{-z^2/2}\) 所定义的分布记为
\(N(0,1)\)，称标准正态。其特征函数为
\(\ee^{-t^2/2}\)：对积分
\(I(t)=\int\ee^{\ii tz}f(z)\,\dd z\) 求导，再用
\(f'(z)=-zf(z)\) 分部积分，得到
\(I'(t)=-tI(t)\)、\(I(0)=1\)，解这条常微分方程即得。若 \(Z\sim N(0,1)\)，则 \(\mu+\sigma Z\) 的分布记为
\(N(\mu,\sigma^2)\)，其中 \(\sigma>0\)。

这里还需要说明：为什么特征函数的极限能够决定分布的极限？对下面标准化后的和，其二阶矩恒为一，因此
\(P(|Y_n|>R)\le R^{-2}\)。这使概率质量不会逃向无穷远，称为\emph{紧性}。任取一个子列，在每个有理数点把分布函数对角抽子列，再利用其单调性，就得到一个进一步子列，在某个分布函数的连续点收敛；紧性保证极限的总质量仍为一。沿这条子列，\(\ee^{\ii tx}\) 是有界连续函数，故特征函数也收敛到该极限分布的特征函数。

最后需用到特征函数的唯一性，它可由一个普通的 Fourier 反演说明。把两个待比较的概率测度分别与方差为 \(\epsilon^2\) 的高斯密度卷积；卷积后的光滑密度的 Fourier 变换是原特征函数乘 \(\ee^{-\epsilon^2t^2/2}\)。这个乘积可积，Fourier 反演表明相同特征函数给出相同卷积密度。令 \(\epsilon\downarrow0\)，高斯核成为近似单位，便得到原概率测度相同。因此，如果一列紧的分布的特征函数趋于某个已知分布的特征函数，所有可能的子列极限都是该分布，整列也就依分布收敛。这是下面所用的特征函数连续性定理在当前情形的证明；其中“紧”不能从任意点态函数极限凭空省去。

\begin{theorem}[独立同分布中心极限定理]
若 \(X_i\) 独立同分布，\(\mathbb E X_i=\mu\)、
\(0<\operatorname{Var}(X_i)=\sigma^2<\infty\)，则
\[
\frac{\sqrt n(\overline X_n-\mu)}{\sigma}\Rightarrow N(0,1).
\]
这是依分布结论，不能读成有限 \(n\) 时统计量精确正态。
\end{theorem}
\begin{proof}
令 \(Y_i=(X_i-\mu)/\sigma\)。其特征函数
\(\varphi(t)=\mathbb E e^{\ii tY_i}\) 在零点满足
\(\varphi(t)=1-t^2/2+o(t^2)\)；这里由 \(\mathbb E Y_i=0\)、
\(\mathbb E Y_i^2=1\) 和受控收敛得到二阶展开，有限二阶矩正用在此处。
独立性给标准化和的特征函数
\[
\varphi(t/\sqrt n)^n
=\bigl(1-t^2/(2n)+o(n^{-1})\bigr)^n
\longrightarrow e^{-t^2/2}.
\]
标准化和的二阶矩为一，故满足前段证明所需的紧性；右侧是标准正态的特征函数。由刚才的子列与唯一性论证，得到所述依分布收敛。
\end{proof}

大数律控制平均值的误差是否消失；中心极限定理描述误差乘以 \(\sqrt n\) 后的极限形状。下节若额外假设原始样本正态，则会得到不必令 \(n\to\infty\) 的精确分布。

\begin{exercise}
证明 \(X_n=n\mathbf1_{A_n}\)、\(P(A_n)=n^{-2}\) 在 \(L^1\) 中收敛到零，却不在 \(L^2\) 中收敛。指出上述中心极限定理在哪一步使用了独立性。
\end{exercise}
